% ============ §1 Ideals (原书 pp.1-6) ============
\section{Ideals}\label{sec:1}\index{Ideals, section 1}

If $A$ is a ring and $I$ an ideal of $A$, it is often important to consider
the residue class ring $A/I$. Set $\bar{A} = A/I$, and write $f\colon A \to \bar{A}$ for the
natural map; then ideals $\bar{J}$ of $\bar{A}$ and ideals $J = f^{-1}(\bar{J})$ of $A$ containing $I$ are
in one-to-one correspondence, with $\bar{J} = J/I$ and $\bar{A}/\bar{J} \cong A/J$. Hence, when
we just want to think about ideals of $A$ containing $I$, it is convenient to
shift attention to $A/I$. (If $I'$ is any ideal of $A$ then $f(I')$ is an ideal of $\bar{A}$,
with $f^{-1}(f(I')) = I + I'$, and $\overline{f(I')} = (I + I')/I$.)

$A$ is itself an ideal of $A$, often written $(1)$ since it is generated by the
identity element $1$. An ideal distinct from $(1)$ is called a \emph{proper ideal}.\index{proper ideal}
An element $a \in A$ which has an inverse in $A$ (that is, for which there exists
$a' \in A$ with $aa' = 1$) is called a \emph{unit}\index{unit} (or \emph{invertible element}) of $A$; this holds
if and only if the principal ideal $(a)$ is equal to $(1)$. If $a$ is a unit and $x$ is
nilpotent then $a + x$ is again a unit: indeed, if $x^n = 0$ then setting $y = -a^{-1}x$,
we have $y^n = 0$; now
\[
  (1 - y)(1 + y + \cdots + y^{n-1}) = 1 - y^n = 1,
\]
so that $a + x = a(1 - y)$ has an inverse.

In a ring $A$ we are allowed to have $1 = 0$, but if this happens then it
follows that $a = 1 \cdot a = 0 \cdot a = 0$ for every $a \in A$, so that $A$ has only one
element $0$; in this case we write $A = 0$. In definitions and theorems about
rings, it may sometimes happen that the condition $A \neq 0$ is omitted even
when it is actually necessary. A ring $A$ is an \emph{integral domain}\index{integral domain} (or simply a
\emph{domain}) if $A \neq 0$, and if $A$ has no zero-divisors other than $0$. If $A$ is an
integral domain and every non-zero element of $A$ is a unit then $A$ is a
\emph{field}.\index{field} A field is characterised by the fact that it is a ring having exactly
two ideals $(0)$ and $(1)$.

An ideal which is maximal among all proper ideals is called a \emph{maximal ideal};\index{maximal ideal}
an ideal $m$ of $A$ is maximal if and only if $A/m$ is a field. Given a
proper ideal $I$, let $\mathcal{M}$ be the set of ideals containing $I$ and not containing
$1$, ordered by inclusion; then Zorn's lemma can be applied to $\mathcal{M}$. Indeed,
$I \in \mathcal{M}$ so that $\mathcal{M}$ is non-empty, and if $\mathcal{L} \subset \mathcal{M}$ is a totally ordered subset
then the union of all the ideals belonging to $\mathcal{L}$ is an ideal of $A$ and obviously
belongs to $\mathcal{M}$, so is the least upper bound of $\mathcal{L}$ in $\mathcal{M}$. Thus by Zorn's
lemma $\mathcal{M}$ has got a maximal element. This proves the following theorem.

\begin{theorem}\label{thm:1.1}
If $I$ is a proper ideal then there exists at least one maximal ideal
containing $I$.
\end{theorem}

An ideal $P$ of $A$ for which $A/P$ is an integral domain is called a \emph{prime ideal}.\index{prime ideal}
In other words, $P$ is prime if it satisfies
\[
  \text{(i) } P \neq A \quad \text{and} \quad \text{(ii) } x, y \notin P \Rightarrow xy \notin P \text{ for } x, y \in A.
\]
A field is an integral domain, so that a maximal ideal is prime.

If $I$ and $J$ are ideals and $P$ a prime ideal, then
\[
  IJ \not\subset P \iff I \not\subset P \text{ or } J \not\subset P.
\]
Indeed, taking $x \in I$ and $y \in J$ with $x, y \notin P$, we have $xy \in IJ$ but $xy \notin P$.

A subset $S$ of $A$ is \emph{multiplicative}\index{multiplicative set} if it satisfies
\[
  \text{(i) } x, y \in S \Rightarrow xy \in S, \quad \text{and} \quad \text{(ii) } 1 \in S;
\]
(here condition (ii) is not crucial: given a subset $S$ satisfying (i), there will
usually not be any essential change on replacing $S$ by $S \cup \{1\}$). If $I$ is
an ideal disjoint from $S$, then exactly as in the proof of Theorem 1 we see
that the set of ideals containing $I$ and disjoint from $S$ has a maximal
element. If $P$ is an ideal which is maximal among ideals disjoint from $S$
then $P$ is prime. For if $x \notin P$, $y \notin P$, then since $P + xA$ and $P + yA$ both
meet $S$, the product $(P + xA)(P + yA)$ also meets $S$. However,
\[
  (P + xA)(P + yA) \subset P + xyA,
\]
so that we must have $xy \notin P$. We have thus obtained the following theorem.

\begin{theorem}\label{thm:1.2}
Let $S$ be a multiplicative set and $I$ an ideal disjoint from $S$;
then there exists a prime ideal containing $I$ and disjoint from $S$.
\end{theorem}

If $I$ is an ideal of $A$ then the set of elements of $A$, some power of which
belongs to $I$, is an ideal of $A$ (for $x^n \in I$ and $y^m \in I \Rightarrow (x + y)^{n+m-1} \in I$
and $(ax)^n \in I$). This set is called the \emph{radical}\index{radical} of $I$, and is sometimes written $\sqrt{I}$:
\[
  \sqrt{I} = \{a \in A \mid a^n \in I \text{ for some } n > 0\}.
\]
If $P$ is a prime ideal containing $I$ then $x^n \in I \subset P$ implies that $x \in P$, and
hence $\sqrt{I} \subset P$; conversely, if $x \notin \sqrt{I}$ then $S = \{1, x, x^2, \ldots\}$ is a multi-
plicative set disjoint from $I$, and by the previous theorem there exists a
prime ideal containing $I$ and not containing $x$. Thus, the radical of $I$ is
the intersection of all prime ideals containing $I$:
\[
  \sqrt{I} = \bigcap_{P \supset I} P.
\]
In particular if we take $I = (0)$ then $\sqrt{(0)}$ is the set of all nilpotent elements
of $A$, and is called the \emph{nilradical}\index{nilradical} of $A$; we will write $\nil(A)$ for this. Then
$\nil(A)$ is the intersection of all the prime ideals of $A$. When $\nil(A) = 0$ we say
that $A$ is \emph{reduced}.\index{reduced ring} For any ring $A$ we write $A_{red}$ for $A/\nil(A)$; $A_{red}$ is of
course reduced.

The intersection of all maximal ideals of a ring $A$ ($\neq 0$) is called the
\emph{Jacobson radical}\index{Jacobson radical}, or simply the radical of $A$, and written $\rad(A)$. If
$x \in \rad(A)$ then for any $a \in A$, $1 + ax$ is an element of $A$ not contained in
any maximal ideal, and is therefore a unit of $A$ by Theorem 1. Conversely
if $x \in A$ has the property that $1 + Ax$ consists entirely of units of $A$ then
$x \in \rad(A)$ (prove this!).

A ring having just one maximal ideal is called a \emph{local ring}\index{local ring}, and a
(non-zero) ring having only finitely many maximal ideals a \emph{semilocal ring}.\index{semilocal ring}
We often express the fact that $A$ is a local ring with maximal ideal $m$ by
saying that $(A, m)$ is a local ring; if this happens then the field $k = A/m$ is
called the \emph{residue field}\index{residue field} of $A$. We will say that $(A, m, k)$ is a local ring to
mean that $A$ is a local ring, $m = \rad(A)$ and $k = A/m$. If $(A, m)$ is a local
ring then the elements of $A$ not contained in $m$ are units; conversely a
(non-zero) ring $A$ whose non-units form an ideal is a local ring.

In general the product $II'$ of two ideals $I$, $I'$ is contained in $I \cap I'$, but
does not necessarily coincide with it. However, if $I + I' = (1)$ (in which case
we say that $I$ and $I'$ are \emph{coprime})\index{coprime ideals}, then $II' = I \cap I'$; indeed, then
\[
  I \cap I' = (I \cap I')(I + I') \subset II' \subset I \cap I'.
\]
Moreover, if $I$ and $I'$, as well as $I$ and $I''$, are coprime, then $I$ and $I'I''$ are coprime:
\[
  (1) = (I + I')(I + I'') \subset I + I'I'' \subset (1).
\]
By induction we obtain the following theorem.

\begin{theorem}\label{thm:1.3}
If $I_1, I_2, \ldots, I_n$ are ideals which are coprime in pairs then
\[
  I_1I_2 \cdots I_n = I_1 \cap I_2 \cap \cdots \cap I_n.
\]
In particular if $A$ is a semilocal ring and $m_1, \ldots, m_n$ are all of its maximal
ideals then
\[
  m_1 \cdots m_n = m_1 \cap \cdots \cap m_n = \rad(A).
\]
\end{theorem}

Furthermore, if $I + I' = (1)$ then
\[
  A/II' \cong A/I \times A/I'.
\]
To see this it is enough to prove that the natural injective map from $A/II' = A/(I \cap I')$ to
$A/I \times A/I'$ is surjective; taking $e \in I$, $e' \in I'$ such that $e + e' = 1$, we
have $ae' + a'e \equiv a \pmod I$, $ae' + a'e \equiv a' \pmod{I'}$ for any $a, a' \in A$,
giving the surjectivity. By induction we get the following theorem.

\begin{theorem}\label{thm:1.4}
If $I_1, \ldots, I_n$ are ideals which are coprime in pairs then
\[
  A/I_1 \cdots I_n \cong A/I_1 \times \cdots \times A/I_n.
\]
\end{theorem}

\begin{example}\label{ex:1.1}
Let $A$ be a ring, and consider the ring $A[[X]]$ of formal power
series over $A$. A power series $f = a_0 + a_1X + a_2X^2 + \cdots$ with $a_i \in A$ is a
unit of $A[[X]]$ if and only if $a_0$ is a unit of $A$. Indeed, if there exists an
inverse $f^{-1} = b_0 + b_1X + \cdots$ then $a_0b_0 = 1$; and conversely if $a_0$ is a unit of $A$, then
\[
  ff^{-1} = a_0b_0 + (a_0b_1 + a_1b_0)X + (a_0b_2 + a_1b_1 + a_2b_0)X^2 + \cdots = 1
\]
can be solved for $b_0, b_1, \ldots$: we just find $b_0, b_1, \ldots$ successively from
$a_0b_0 = 1$, $a_0b_1 + a_1b_0 = 0, \ldots$

Since the formal power series ring in several variables $A[[X_1, \ldots, X_n]]$ can be thought of as
$(A[[X_1, \ldots, X_{n-1}]])[[X_n]]$, here also $f = a_0 + \sum a_iX_i + \sum a_{ij}X_iX_j + \cdots$ is a unit
if and only if the constant term $a_0$ is a unit of $A$; from this we see that
if $g \in (X_1, \ldots, X_n)$ then $1 + gh$ is a unit for any power series $h$, so that
$g \in \rad(A[[X_1, \ldots, X_n]])$, and hence
\[
  (X_1, \ldots, X_n) \subset \rad(A[[X_1, \ldots, X_n]]).
\]
If $k$ is a field then $k[[X_1, \ldots, X_n]]$ is a local ring with maximal ideal
$(X_1, \ldots, X_n)$. If $A$ is any ring and we set $B = A[[X_1, \ldots, X_n]]$, then since
any maximal ideal of $B$ contains $(X_1, \ldots, X_n)$, it corresponds to a maximal
ideal of $B/(X_1, \ldots, X_n) \cong A$, and so is of the form $mB + (X_1, \ldots, X_n)$, where
$m$ is a maximal ideal of $A$. If we write $m$ for this then $m \cap A = m$.

By contrast the case of polynomial rings is quite complicated; here it
is just not true that a maximal ideal of $A[X]$ must contain $X$. For example,
$X - 1$ is a non-unit of $A[X]$, and so there exists a maximal ideal $m$
containing it, and $X \notin m$. Also, if $m$ is a maximal ideal of $A[X]$, it does
not necessarily follow that $m \cap A$ is a maximal ideal of $A$.

If $A$ is an integral domain then so are both $A[X]$ and $A[[X]]$: if
$f = a_rX^r + a_{r+1}X^{r+1} + \cdots$ and $g = b_sX^s + b_{s+1}X^{s+1} + \cdots$ with $a_r \neq 0$,
$b_s \neq 0$ then $fg = a_rb_sX^{r+s} + \cdots \neq 0$. If $I$ is an ideal of $A$ we write $I[X]$ or
$I[[X]]$ for the set of polynomials or power series with coefficients in $I$;
these are ideals of $A[X]$ or $A[[X]]$, the kernels of the homomorphisms
\[
  A[X] \to (A/I)[X] \quad \text{or} \quad A[[X]] \to (A/I)[[X]]
\]
obtained by reducing coefficients modulo $I$. Hence
\[
  \sqrt{I[X]} = (\sqrt{I})[X], \quad \text{and} \quad A[X]/I[X] \cong (A/I)[X];
\]
in particular if $P$ is a prime ideal then $P[X]$ and $P[[X]]$ are prime ideals
of $A[X]$ and $A[[X]]$, respectively.

If $I$ is finitely generated, that is $I = a_1A + \cdots + a_lA$, then $I[[X]] =
a_1A[[X]] + \cdots + a_lA[[X]] = I \cdot A[[X]]$; however, if $I$ is not finitely generated
then $I[[X]]$ is bigger than $I \cdot A[[X]]$. In the polynomial ring this distinction
does not arise, and we always have $I[X] = I \cdot A[X]$.
\end{example}

\begin{example}\label{ex:1.2}
For a ring $A$ and $a, b \in A$, we have $aA \subset bA$ if and only if $a$
is divisible by $b$, that is $a = bc$ for some $c \in A$. We assume that $A$ is an
integral domain in what follows. An element $a \in A$ is said to be \emph{irreducible}\index{irreducible element}
if $a$ is not a unit of $A$ and satisfies the condition
\[
  a = bc \Rightarrow b \text{ or } c \text{ is a unit of } A.
\]
This is equivalent to saying that $aA$ is maximal among proper principal ideals. If $aA$ is a
prime ideal then $a$ is said to be \emph{prime}.\index{prime element} As one sees easily,
a prime element is irreducible, but the converse does not always hold.

Suppose that an element $a$ has two expressions as products of prime elements:
\[
  a = p_1p_2 \cdots p_n = p_1' \cdots p_m', \quad \text{with } p_i \text{ and } p_i' \text{ prime}.
\]
Then $n = m$, and after a suitable reordering of the $p_i'$ we have $p_iA = p_i'A$;
for $p_1 \cdots p_n$ is divisible by $p_1'$, and so one of the factors, say $p_i$, is
divisible by $p_1'$. Now since both $p_1'$ and $p_i$ are irreducible, $p_1'A = p_iA$,
hence $p_i = up_1'$, with $u$ a unit, and $p_2 \cdots p_n = up_2' \cdots p_m'$. We can replace
$p_i$ by $up_i$, and induction on $n$ completes the proof. In this sense,
factorisation into prime elements (whenever possible) is unique.

An integral domain in which any element which is neither $0$ nor a unit
can be expressed as a product of prime elements is called a \emph{unique factorisation
domain}\index{unique factorisation domain} (abbreviated to UFD), or a \emph{factorial ring}. It is well
known that a principal ideal domain, that is an integral domain in which
every ideal is principal, is a UFD (see Ex.~\ref{ex:1.4-ex}). If $A$ is a principal ideal
domain then the prime ideals are of the form $(0)$ or $pA$ with $p$ a prime
element, and the latter are maximal ideals.

If $k$ is a field then $k[X_1, \ldots, X_n]$ is a UFD, as is well-known (see Ex. 20.2).
If $f(X_1, \ldots, X_n)$ is an irreducible polynomial then $(f)$ is a prime ideal,
but is not maximal if $n > 1$ (see \S\ref{sec:5}).

$\Z[\sqrt{-5}]$ is not a UFD; indeed if $z = n + m\sqrt{-5}$ with $n, m \in \Z$ then
$z\bar{z} = n^2 + 5m^2$, and since $2 = n^2 + 5m^2$ has no integer solutions it
follows that $2$ is an irreducible element of $\Z[\sqrt{-5}]$, but we see from
$2 \cdot 3 = (1 + \sqrt{-5})(1 - \sqrt{-5})$ that $2$ is not a prime element. We write
\[
  A = \Z[\sqrt{-5}] = \Z[X]/(X^2 + 5);
\]
then setting $k = \Z/2\Z$ we have
\[
  A/2A = \Z[X]/(2, X^2 + 5) = k[X]/(X^2 - 1) = k[X]/(X - 1)^2.
\]
Then $P = (2, 1 - \sqrt{-5})$ is a maximal ideal of $A$ containing $2$.
\end{example}

% ============ 习题 ============
\section*{Exercises to \S\ref{sec:1}}
\emph{Prove the following propositions.}

\begin{exercise}\label{ex:1.1-ex}
Let $A$ be a ring, and $I \subset \nil(A)$ an ideal made up of nilpotent elements;
if $a \in A$ maps to a unit of $A/I$ then $a$ is a unit of $A$.
\end{exercise}

\begin{exercise}\label{ex:1.2-ex}
Let $A_1, \ldots, A_n$ be rings; then the prime ideals of $A_1 \times \cdots \times A_n$ are of the form
\[
  A_1 \times \cdots \times A_{i-1} \times P_i \times A_{i+1} \times \cdots \times A_n,
\]
where $P_i$ is a prime ideal of $A_i$.
\end{exercise}

\begin{exercise}\label{ex:1.3-ex}
Let $A$ and $B$ be rings, and $f\colon A \to B$ a surjective homomorphism.
\begin{enumerate}[(a)]
  \item Prove that $f(\rad A) \subset \rad B$, and construct an example where the
        inclusion is strict.
  \item Prove that if $A$ is a semilocal ring then $f(\rad A) = \rad B$.
\end{enumerate}
\end{exercise}

\begin{exercise}\label{ex:1.4-ex}
Let $A$ be an integral domain. Then $A$ is a UFD if and only if every
irreducible element is prime and the principal ideals of $A$ satisfy the
ascending chain condition. (Equivalently, every non-empty family of
principal ideals has a maximal element.)
\end{exercise}

\begin{exercise}\label{ex:1.5-ex}
Let $\{P_\lambda\}_{\lambda \in \Lambda}$ be a non-empty family of prime ideals, and suppose that the $P_\lambda$
are totally ordered by inclusion; then $\bigcap P_\lambda$ is a prime ideal. Also, if $I$ is
any proper ideal, the set of prime ideals containing $I$ has a minimal element.
\end{exercise}

\begin{exercise}\label{ex:1.6-ex}
Let $A$ be a ring, $I, P_1, \ldots, P_n$ ideals of $A$, and suppose that $P_1, \ldots, P_n$ are
prime, and that $I$ is not contained in any of the $P_i$; then there exists an
element $x \in I$ not contained in any $P_i$.
\end{exercise}
